\chapter{Canker Sores (Aphthous Ulcers)}

\section{Definition}
Canker sores (aphthous ulcers) are shallow, painful ulcers that develop on the oral mucosa. They affect approximately 20 percent of the population and are classified as minor (less than 1 cm), major (greater than 1 cm), or herpetiform.

\section{Pathophysiology}
The exact pathogenesis involves an immune-mediated attack on oral epithelial cells, driven by T-cell activation and cytokine release. Focal necrosis of the epithelium results in a painful, crater-like defect covered by a grayish fibrinous membrane surrounded by an erythematous halo. The ulcers heal spontaneously by epithelial regeneration over one to three weeks.

\section{Etiology}
\begin{itemize}
\item Minor trauma (biting the cheek, sharp foods, dental appliances)
\item Stress and anxiety
\item Hormonal changes (menstruation, pregnancy)
\item Food sensitivities (spicy, acidic foods; chocolate, nuts, cheese)
\item Nutritional deficiencies (vitamin B12, folate, iron, zinc)
\item Gluten sensitivity/celiac disease
\item Immunodeficiency states
\item Medications (NSAIDs, beta-blockers)
\item Family history (genetic predisposition)
\item Autoimmune conditions (Behcet disease, IBD)
\end{itemize}

\section{Indications for Evaluation}
\begin{itemize}
\item Painful, shallow ulcers on non-keratinized oral mucosa
\item Ulcers with a grayish-white center and red border
\item Recurrent episodes lasting 7-14 days
\item Difficulty eating, drinking, or speaking due to pain
\item Prodromal burning or tingling sensation
\item Ulcers with fever, rash, or genital ulcers (possible Behcet disease)
\item Ulcers persisting longer than two weeks
\end{itemize}

\section{Related Symptoms}
\begin{itemize}
\item Oral pain (mouth pain)
\item Difficulty swallowing (dysphagia)
\item Mouth sores (oral ulcerations)
\item Cold sores (herpes labialis)
\item Glossitis (tongue inflammation)
\item Gingivitis (gum inflammation)
\item Behcet disease
\end{itemize}

\section{Self-Care/Home Management}
\begin{itemize}
\item Avoid spicy, acidic, or rough foods during outbreaks
\item Use over-the-counter topical preparations (benzocaine)
\item Apply milk of magnesia or baking soda paste
\item Rinse with salt water several times daily
\item Use a soft-bristled toothbrush
\item Identify and avoid personal dietary triggers
\item Consider vitamin B12, folate, or iron supplements if deficient
\end{itemize}

\section{Diagnostic Approach}
\begin{itemize}
\item Clinical diagnosis based on appearance and history
\item Biopsy for atypical, persistent, or unusually large ulcers
\item Blood tests for nutritional deficiencies
\item Tests for celiac disease if gastrointestinal symptoms present
\item Evaluation for Behcet disease if oral and genital ulcers coexist
\end{itemize}

\section{Treatment/Management Overview}
\begin{itemize}
\item Topical corticosteroids (triamcinolone acetonide)
\item Topical analgesics (lidocaine, benzocaine)
\item Oral corticosteroid short course for severe cases
\item Immunomodulatory agents (colchicine) for refractory cases
\item Antimicrobial mouth rinses (chlorhexidine)
\item Chemical cautery (silver nitrate) for rapid pain relief
\item Nutritional supplementation if deficiencies identified
\end{itemize}
